% Receptive-field figure: two stacked 3x3 conv layers on the 11x11 toy grid.
% Bottom (input) cells are filled with PAN values from
% analysis/output/metrics-diagram-pan-values.tex.
\begin{tikzpicture}[
    every node/.style={font=\footnotesize},
    x=0.50cm, y=0.50cm,
]
    \def\skewx{0.45}
    \def\skewy{0.40}
    \def\layersep{6.5}

    \foreach \L in {0,1,2} {
        \foreach \x in {-1,...,11} {
            \foreach \y in {-1,...,11} {
                \pgfmathsetmacro{\X}{\x + \skewx*\y}
                \pgfmathsetmacro{\Y}{\skewy*\y + \layersep*\L}
                \coordinate (n\L-\x-\y) at (\X, \Y);
            }
        }
    }

    \newcommand{\panvalue}[3]{%
        \pgfmathtruncatemacro{\xm}{#1-1}%
        \pgfmathtruncatemacro{\ym}{#2-1}%
        \filldraw[draw=black!50, line width=0.15pt, fill=white!#3!black]
            (n0-\xm-\ym) -- (n0-#1-\ym) -- (n0-#1-#2) -- (n0-\xm-#2) -- cycle;
    }
    \input{analysis/output/metrics-diagram-pan-values.tex}

    \foreach \L in {1, 2} {
        \foreach \x in {1,...,10} {
            \foreach \y in {1,...,10} {
                \pgfmathtruncatemacro{\xm}{\x-1}
                \pgfmathtruncatemacro{\ym}{\y-1}
                \filldraw[draw=black!55, line width=0.15pt, fill=white]
                    (n\L-\xm-\ym) -- (n\L-\x-\ym) -- (n\L-\x-\y) -- (n\L-\xm-\y) -- cycle;
            }
        }
    }

    \def\ax{6}
    \def\ay{6}
    \pgfmathtruncatemacro{\axm}{\ax-1}
    \pgfmathtruncatemacro{\aym}{\ay-1}
    \pgfmathtruncatemacro{\plo}{\ax-2}
    \pgfmathtruncatemacro{\phi}{\ax+1}
    \pgfmathtruncatemacro{\rlo}{\ax-3}
    \pgfmathtruncatemacro{\rhi}{\ax+2}

    \filldraw[draw=red!70!black, line width=0.4pt, fill=red!60]
        (n1-\axm-\aym) -- (n1-\ax-\aym) -- (n1-\ax-\ay) -- (n1-\axm-\ay) -- cycle;

    \filldraw[draw=blue!70!black, line width=0.4pt, fill=blue!60]
        (n2-\axm-\aym) -- (n2-\ax-\aym) -- (n2-\ax-\ay) -- (n2-\axm-\ay) -- cycle;

    \draw[red!75!black, line width=0.9pt]
        (n0-\plo-\plo) -- (n0-\phi-\plo) -- (n0-\phi-\phi) -- (n0-\plo-\phi) -- cycle;

    \draw[blue!75!black, line width=0.9pt]
        (n1-\plo-\plo) -- (n1-\phi-\plo) -- (n1-\phi-\phi) -- (n1-\plo-\phi) -- cycle;

    \draw[orange!90!black, line width=1.1pt]
        (n0-\rlo-\rlo) -- (n0-\rhi-\rlo) -- (n0-\rhi-\rhi) -- (n0-\rlo-\rhi) -- cycle;

    \coordinate (apex1) at ($(n1-\axm-\aym)!0.5!(n1-\ax-\ay)$);
    \coordinate (apex2) at ($(n2-\axm-\aym)!0.5!(n2-\ax-\ay)$);

    \draw[red!75!black, line width=0.4pt] (apex1) -- (n0-\plo-\plo);
    \draw[red!75!black, line width=0.4pt] (apex1) -- (n0-\phi-\plo);
    \draw[red!75!black, line width=0.4pt] (apex1) -- (n0-\phi-\phi);
    \draw[red!75!black, line width=0.4pt] (apex1) -- (n0-\plo-\phi);

    \draw[blue!75!black, line width=0.4pt] (apex2) -- (n1-\plo-\plo);
    \draw[blue!75!black, line width=0.4pt] (apex2) -- (n1-\phi-\plo);
    \draw[blue!75!black, line width=0.4pt] (apex2) -- (n1-\phi-\phi);
    \draw[blue!75!black, line width=0.4pt] (apex2) -- (n1-\plo-\phi);

    \node[anchor=east] at (n0--1-5) {Input};
    \node[anchor=east] at (n1--1-5) {Layer 1};
    \node[anchor=east] at (n2--1-5) {Layer 2};

    \node[anchor=west, red!70!black] at ($(n0-\rhi-\rhi)+(0.8,0.5)$)
        {$3 \times 3$ patch};
    \node[anchor=west, orange!80!black, align=left] at ($(n0-\rhi-\rlo)+(0.8,-0.6)$)
        {$5 \times 5$ effective\\input window};

\end{tikzpicture}
